% sec-03: the thirty-day pipeline.  Figure 15 (d2-type layers), Table 22.
\section{The Thirty-Day Pipeline}

One rule governs this section: an item whose next action belongs to this company
is not pipeline. It is work, and it belongs on a to-do list. Confusing the two
produces a pipeline that looks busy and moves nothing, and a founder who reads it
feels productive without being so.

Every row below therefore waits on somebody else.

\begin{apptable}
\begin{tabularx}{\textwidth}{>{\raggedright\arraybackslash}p{3.4cm} >{\raggedright\arraybackslash}p{1.6cm} >{\raggedright\arraybackslash}p{2.6cm} >{\raggedright\arraybackslash}X}
\headrow \thead{Stage} & \thead{Count} & \thead{Interval basis} & \thead{What the stage contains}\\
\midrule
Asked and waiting & 17 & Varies by counterparty & Every open item across federal, commercial, institutional, state and regional \\
Answerable inside thirty days & 9 & 5 to 15 business days & Those whose stated interval is short enough to resolve in the window \\
Unblocks something else & 4 & No interval; each is routed & The four diligence questions with no answer today \\
Decisions expected & 5 & If nothing goes wrong & One classification, one determination, one instrument, one site routing, and zero to three federal replies \\
Cannot be chased at all & 4 & Somebody else's calendar & Classification, foundation cycle, developer window, review board calendar \\
\end{tabularx}
\end{apptable}
\vspace{-0.60cm}
\tabcap{Table 22. Five pipeline stages with their counts and the basis of each\\
interval, including the four items no action by this company can advance.}

\begin{appfloat}[!tb]
\begin{appfig}[d2-type / layered stack with a count on each layer]
% Four stacked layers of decreasing width on a common axis, and a fifth drawn to
% one side because it is not part of the funnel.  There is no edge to the side
% layer: the absence of the connector is the point.
\node[d2step,text width=88mm,align=center] (l1) at (0,0) {Asked and waiting};
\node[d2step,text width=68mm,align=center] (l2) at (0,-1.15) {Answerable inside thirty days};
\node[d2step,text width=48mm,align=center] (l3) at (0,-2.30) {Unblocks something else};
\node[d2key,text width=34mm,align=center]  (l4) at (0,-3.45) {Decisions expected};
\node[d2ghost,text width=42mm,align=center] (l5) at (5.10,-1.72) {Cannot be chased at all};
\node[d2cellk,minimum width=12mm,text width=10mm] at (5.05,0) {17};
\node[d2cellk,minimum width=12mm,text width=10mm] at (4.35,-1.15) {9};
\node[d2cellk,minimum width=12mm,text width=10mm] at (3.65,-2.30) {4};
\node[d2cellk,minimum width=12mm,text width=10mm] at (2.95,-3.45) {5};
\node[d2cellg,minimum width=12mm,text width=10mm] at (7.35,-1.72) {4};
\draw[d2edgeb] (l1) -- (l2);
\draw[d2edgeb] (l2) -- (l3);
\draw[d2edgeb] (l3) -- (l4);
\node[font=\tiny\itshape,text=fundgray,anchor=north west,text width=118mm]
  at (-4.60,-4.15) {No edge reaches the side layer, and the absence is the point.
  The four items there are not a stage of the funnel: they arrive on somebody
  else's calendar, and drawing a connector would suggest they can be moved along.};
\ptitle{-4.60}{0.72}{Wide where questions are cheap, narrow where decisions are dear}
\end{appfig}
\vspace{-0.60cm}
\figcaption{Figure 15. The thirty-day pipeline as four narrowing layers, with the four\\
items that cannot be chased drawn outside the funnel rather than inside it.}
\end{appfloat}

\subsection{The four that cannot be chased, named separately}

The FDA classification arrives when it arrives, and an applicant who chases a
classification request is an applicant the office remembers
\cite{fdacombination2026}. A foundation cycle runs on published dates that are
not negotiable. A developer submission window opens on the developer's calendar.
An institutional review board meets when it meets.

Naming these four separately matters for a practical reason. A founder with time
and anxiety will chase something, and it should not be one of these.

\subsection{What thirty days produce if nothing goes wrong}

Five decisions, four of which unblock something else. That is a reasonable month
for a company of one person. If nothing arrives at all, the four unanswered
diligence questions are unchanged and the stop conditions are unchanged, which is
the honest downside and is already written down in the preceding day's packet
rather than discovered later.

\subsection{Why the pipeline is published rather than kept internally}

A pipeline in a public repository is unusual and the reason is deliberate. A
counterparty who reads it can see exactly where their own item sits, what
interval this company has assigned to it, and that no action is planned before
that interval elapses. That removes the most common source of friction in early
correspondence, which is a recipient wondering whether silence will be read as a
no and answered with a chase.

It also constrains this company. An interval published against a named
counterparty is an interval that cannot quietly shorten because a founder became
impatient on a Tuesday, and the four unchaseable items are protected by exactly
that constraint.

